\documentclass[11pt]{article}
\usepackage[margin=1in]{geometry}
\usepackage[T1]{fontenc}
\usepackage[utf8]{inputenc}
\usepackage{lmodern}
\usepackage{longtable}
\usepackage{booktabs}
\usepackage{enumitem}
\usepackage{hyperref}
\hypersetup{colorlinks=true,linkcolor=blue,urlcolor=blue}

\title{AtlasAutoware System Guide for Beginners\\
(How the Self-Driving Racecar Works and What Every Code File Does)}
\author{AtlasAutoware Repository Documentation}
\date{\today}

\begin{document}
\maketitle
\tableofcontents
\newpage

\section{Purpose of this document}
This guide explains the AtlasAutoware codebase for someone new to self-driving cars.
It covers:
\begin{itemize}
  \item core autonomous-driving vocabulary,
  \item how the end-to-end racing stack works,
  \item what runs in simulation vs. real hardware,
  \item the main control loop that runs on the vehicle,
  \item and what each code file is responsible for.
\end{itemize}

\section{Beginner overview: how a self-driving racecar works}
A self-driving car continuously repeats a control cycle:
\begin{enumerate}
  \item \textbf{Sense}: read lidar, camera, IMU, and odometry.
  \item \textbf{Estimate state}: compute where the car is, how fast it is moving, and how it is rotating.
  \item \textbf{Plan}: select a path and target speed (here, a precomputed raceline with safety logic).
  \item \textbf{Control}: compute steering and speed commands.
  \item \textbf{Actuate}: send commands to motor controller and steering.
  \item \textbf{Repeat}: run again at fixed rate (about 50 Hz in this stack).
\end{enumerate}

AtlasAutoware is built for F1TENTH (1/10-scale racing), but it uses the same architecture patterns used in full-size autonomous systems: modular sensing, planning, control, and safety layers.

\section{Vocabulary and core terms}
\begin{description}[leftmargin=2.2cm,style=nextline]
  \item[ROS 2] Middleware that lets independent software nodes exchange messages over topics.
  \item[Node] A program process with one job (for example, lidar driver or controller).
  \item[Topic] Named message channel (example: \texttt{/scan}, \texttt{/drive}).
  \item[Publisher] Node that sends messages to a topic.
  \item[Subscriber] Node that receives messages from a topic.
  \item[Message] Typed data packet (LaserScan, Odometry, AckermannDriveStamped, etc.).
  \item[Frame] Coordinate system name in TF (example: \texttt{map}, \texttt{base\_link}).
  \item[TF] ROS transform tree that relates frames over time.
  \item[Lidar] Laser range sensor measuring distances around the car.
  \item[IMU] Inertial Measurement Unit (accelerometer + gyroscope).
  \item[Odometry] Motion estimate (velocity/pose) from wheel or state estimator.
  \item[Raceline] Fastest practical closed-loop path around a track.
  \item[Curvature] How sharply a path bends; high curvature means tighter turn (denoted $\kappa$).
  \item[Trajectory] Time-parameterized path (position and often speed profile).
  \item[Waypoint] One sampled point along a path/raceline.
  \item[Control loop] Repeating cycle that computes commands from latest sensor/state data.
  \item[MPC] Model Predictive Control; solves an optimization each cycle to track path with constraints.
  \item[MPCC] Model Predictive Contouring Control; MPC variant optimizing contouring/lag errors.
  \item[OSQP] Quadratic-program solver used by MPC in this repository.
  \item[Fallback controller] Backup controller used when primary solver/path fails.
  \item[MAP controller] Model- and Acceleration-based Pursuit; robust geometric+tire-aware fallback.
  \item[Pure Pursuit] Geometric path-following controller using lookahead point.
  \item[Stanley controller] Steering law based on heading + cross-track error.
  \item[AEB] Automatic Emergency Braking; stops if obstacle is too close ahead.
  \item[Latency compensation] Predicting future vehicle state to account for actuation delay.
  \item[Traction governor] Speed scaler that reduces command when measured lateral acceleration exceeds grip budget.
  \item[Occupancy grid] 2D map where cells are free, occupied, or unknown.
  \item[SLAM] Simultaneous Localization and Mapping.
  \item[Loop closure] SLAM correction when revisiting known places.
  \item[Deskew] Correcting rolling lidar scan distortion caused by motion during scan.
  \item[EKF] Extended Kalman Filter; state estimator combining model + noisy sensor measurements.
  \item[Slip] Wheel speed diverges from true ground speed due to tire traction loss.
  \item[Friction limit] Maximum tire force before sliding.
  \item[Friction-limited profile] Speed profile that respects lateral/longitudinal grip limits.
  \item[Disparity extender] Follow-the-gap safety logic inflating obstacle edges by vehicle width.
  \item[YOLO] Real-time object detector used for opponent car perception.
  \item[TensorRT] NVIDIA runtime for optimized neural inference on Jetson GPU.
  \item[VESC] Motor controller for throttle/motor actuation and telemetry.
  \item[PCA9685] PWM board for generating servo/PPM pulses.
\end{description}

\section{System architecture in this repository}
\subsection{Main execution modes}
\begin{itemize}
  \item \textbf{Simulation mode}: \texttt{gym\_bridge} publishes simulated lidar/odom and receives \texttt{/drive}.
  \item \textbf{Real-car mode}: hardware drivers publish real sensors and \texttt{drive\_node} actuates the car.
  \item \textbf{Mapping mode}: \texttt{slam\_toolbox} + \texttt{mapping\_driver} + \texttt{track\_learner} build map/raceline.
  \item \textbf{Racing mode}: \texttt{raceline\_mpc} follows optimized raceline with safety/fallback logic.
\end{itemize}

\subsection{Primary dataflow on real vehicle}
\begin{enumerate}
  \item \texttt{rplidar\_node} publishes \texttt{/scan}.
  \item \texttt{oakd\_camera} publishes \texttt{/oakd/rgb}, \texttt{/oakd/camera\_info}, \texttt{/oakd/imu}.
  \item \texttt{velocity\_ekf} fuses IMU + wheel odometry and publishes \texttt{/ekf/odom}.
  \item Localization (typically particle filter, external) publishes map-frame odometry (configured as \texttt{/pf/pose/odom}).
  \item \texttt{raceline\_mpc} consumes odometry + scan (+ optional IMU), computes \texttt{/drive}.
  \item \texttt{drive\_node} consumes \texttt{/drive}, sends hardware actuation through PCA9685 or direct VESC UART.
\end{enumerate}

\section{The main code that runs the vehicle loop}
The core race-time loop on real car is in \texttt{f1tenth\_gym\_ros/raceline\_mpc.py}:
\begin{itemize}
  \item Node class: \texttt{RacelineMPC}
  \item Timer callback: \texttt{\_loop()} at \texttt{control\_hz} (default 50 Hz)
\end{itemize}

Each cycle, \texttt{\_loop()} does:
\begin{enumerate}
  \item Validate fresh scan and odometry.
  \item Predict current state forward by actuation delay (command-buffer based latency compensation).
  \item Find nearest raceline index.
  \item Try MPC solve (\texttt{KinematicMPC.solve}); if unavailable/fails, use \texttt{MAPController.control} fallback.
  \item Scale speed by \texttt{v\_scale}.
  \item Apply traction governor scaling from IMU yaw-rate if enabled.
  \item Run AEB check in forward cone; force stop if obstacle within dynamic stop distance.
  \item Publish \texttt{AckermannDriveStamped} to \texttt{/drive}.
\end{enumerate}

The actuator endpoint loop is in \texttt{f1tenth\_gym\_ros/drive\_node.py}:
\begin{itemize}
  \item Subscribes to \texttt{/drive}.
  \item Converts speed/steering to hardware commands.
  \item Enforces arming hold, watchdog timeout, and neutral-on-shutdown safety.
\end{itemize}

Together, these two files implement the practical real-car loop: \emph{compute command} then \emph{execute command safely}.

\section{Detailed file-by-file code reference}
\subsection{Core runtime package: \texttt{f1tenth\_gym\_ros/}}
\begin{longtable}{p{0.30\linewidth}p{0.66\linewidth}}
\toprule
\textbf{File} & \textbf{Role and important details} \\
\midrule
\endfirsthead
\toprule
\textbf{File} & \textbf{Role and important details} \\
\midrule
\endhead
\texttt{\_\_init\_\_.py} & Package marker for Python module import. \\
\texttt{camera\_perception.py} & YOLO opponent perception node for real-car operation. Supports TensorRT, CUDA, or CPU inference; converts image detections to relative/world opponent positions; publishes marker and pose topics for race logic. \\
\texttt{drive\_node.py} & Unified actuation node. Auto-detects PCA9685 vs VESC backend, applies throttle/steer mapping, watchdog neutraling, startup arming delay, and optional VESC telemetry publish. \\
\texttt{gym\_bridge.py} & Simulation bridge between f1tenth gym and ROS topics. Publishes scans/odom/TF and accepts drive commands; supports one or two agents. \\
\texttt{map\_controller.py} & Model- and Acceleration-based Pursuit (MAP) fallback controller. Builds lateral-acceleration LUT and inverts desired acceleration to steering with curvature-aware lookahead. \\
\texttt{mapping\_driver.py} & Conservative autonomous mapper using disparity-extender gap-following to produce safe SLAM laps. \\
\texttt{mpc\_controller.py} & Kinematic LTV-MPC implementation, latency prediction utility, and traction governor. Builds persistent OSQP QP for fast repeated solves. \\
\texttt{mpcc\_controller.py} & MPCC implementation (contouring-control variant) for advanced optimization experiments/benchmarks. \\
\texttt{oakd\_camera.py} & OAK-D Pro driver node. Publishes RGB, camera intrinsics, IMU, and optional on-device YOLO detections from Myriad X VPU. \\
\texttt{opponent\_driver.py} & Secondary/opponent vehicle driver for two-car simulation scenarios and race-agent testing. \\
\texttt{pca9685.py} & Low-level PCA9685 PWM driver and conversion helpers from desired speed/steer to pulse widths. Pure logic, unit tested. \\
\texttt{pursuit\_agent.py} & Progressive-learning racing agent and shared raceline utilities (load/save/find nearest/lookahead). Used by other modules for common raceline operations. \\
\texttt{race\_agent.py} & Opponent-aware race policy node with CRUISE/ATTACK/DEFEND/EVADE strategy integration and visualization outputs. \\
\texttt{race\_brain.py} & Pure-logic opponent tracking and strategic decision engine; lidar clustering, track association, fusion with camera, and mode decisions. \\
\texttt{raceline\_mpc.py} & Main competition time-trial node. Loads raceline, optional refinement and reprofile, runs MPC with MAP fallback, traction governor, and lidar AEB, then publishes drive commands. \\
\texttt{raceline\_optimizer.py} & Offline map-to-raceline optimizer: corridor extraction, centerline building, min-curvature optimization, and speed profiling to write raceline CSVs. \\
\texttt{raceline\_refiner.py} & Corridor-bounded minimum-curvature refinement on an existing raceline. \\
\texttt{rplidar\_node.py} & RPLidar ROS driver; bins irregular scan points into fixed-angle LaserScan grid, optional EKF-based deskew, auto-reconnect thread. \\
\texttt{scan\_deskew.py} & Beam-wise geometric deskew utilities for correcting lidar motion distortion during each scan revolution. \\
\texttt{spliner.py} & Frenet-frame overtaking-path deformation helper (``spliner'' strategy). \\
\texttt{track\_learner.py} & Live raceline learner during mapping sessions. Watches \texttt{/map}, periodically runs optimizer, tracks best lap estimate, and promotes best raceline. \\
\texttt{velocity\_ekf.py} & EKF state estimator for body-frame velocities and yaw-rate with slip-aware wheel-speed gating and nonholonomic pseudo-measurement. \\
\texttt{velocity\_profiler.py} & Friction-limited forward-backward velocity profile generator over curvature/path length. \\
\texttt{vesc\_protocol.py} & VESC UART packet encode/decode, command builders, telemetry parsing. Pure logic, hardware-independent. \\
\bottomrule
\end{longtable}

\subsection{Launch files: \texttt{launch/}}
\begin{longtable}{p{0.30\linewidth}p{0.66\linewidth}}
\toprule
\textbf{File} & \textbf{Role and important details} \\
\midrule
\endfirsthead
\toprule
\textbf{File} & \textbf{Role and important details} \\
\midrule
\endhead
\texttt{car\_bringup\_launch.py} & Real-car stack launcher. Starts lidar, camera, drive node, EKF, and raceline MPC with configurable toggles and static sensor transforms. \\
\texttt{gym\_bridge\_launch.py} & Standard simulation launch with gym bridge, map server/lifecycle manager, RViz, and robot state publishers. \\
\texttt{headless\_bridge\_launch.py} & Simulation launch without RViz (useful for noVNC/benchmarking/automation environments). \\
\texttt{slam\_mapping\_launch.py} & Mapping launch for \texttt{slam\_toolbox} plus optional autonomous mapping driver and track learner. \\
\bottomrule
\end{longtable}

\subsection{Configuration files: \texttt{config/}}
\begin{longtable}{p{0.30\linewidth}p{0.66\linewidth}}
\toprule
\textbf{File} & \textbf{Role and important details} \\
\midrule
\endfirsthead
\toprule
\textbf{File} & \textbf{Role and important details} \\
\midrule
\endhead
\texttt{sim.yaml} & Parameters for simulation bridge: map path, start poses, scan settings, topic names, teleop enable, and number of agents. \\
\texttt{hardware.yaml} & Real-car parameters: actuation calibration, hardware ports, sensor topics, EKF, perception, and racing-controller settings (including delay and profiling). \\
\texttt{slam\_mapping.yaml} & \texttt{slam\_toolbox} tuning for online mapping (frames, update intervals, loop closure, scan matching). \\
\bottomrule
\end{longtable}

\subsection{Tools and offline scripts: \texttt{tools/}}
\begin{longtable}{p{0.30\linewidth}p{0.66\linewidth}}
\toprule
\textbf{File} & \textbf{Role and important details} \\
\midrule
\endfirsthead
\toprule
\textbf{File} & \textbf{Role and important details} \\
\midrule
\endhead
\texttt{annotate\_raceline.py} & Draws and labels raceline features (corners/apex/speed annotations) onto map imagery. \\
\texttt{auto\_tune.py} & Automatic tuning loop for racing parameters using repeated benchmark/evaluation runs. \\
\texttt{benchmark\_delay.py} & Measures effect of actuation delay and delay compensation strategies. \\
\texttt{benchmark\_lap.py} & Runs lap-time benchmarking of current raceline/controller configuration in simulation. \\
\texttt{benchmark\_lookahead.py} & Evaluates curvature-aware lookahead behavior in MAP controller. \\
\texttt{benchmark\_map.py} & Benchmarks MAP steering LUT interpolation choices. \\
\texttt{benchmark\_mpcc.py} & Compares MPCC and tracking MPC under matched conditions. \\
\texttt{benchmark\_refiner.py} & Benchmarks raceline refinement pipeline and constraints. \\
\texttt{benchmark\_spliner.py} & Benchmarks spliner-based overtaking strategy effects. \\
\texttt{collect\_camera\_data.py} & Captures camera data for perception model training datasets. \\
\texttt{draw\_raceline.py} & Visualizes raceline CSV over track/map image for inspection. \\
\texttt{image\_to\_map.py} & Converts rendered track drawings into occupancy-map format for optimization/SLAM tooling. \\
\texttt{refine\_comp\_raceline.py} & End-to-end competition raceline cleanup: refine, reprofile, validate, install. \\
\texttt{reprofile\_raceline.py} & Recomputes only raceline speed column via friction-limited profiler. \\
\texttt{train\_car\_detector.py} & Trains/exports YOLO opponent detector model for deployment (ONNX/TensorRT pipeline). \\
\bottomrule
\end{longtable}

\subsection{Tests: \texttt{tests/} and \texttt{test/}}
\begin{longtable}{p{0.30\linewidth}p{0.66\linewidth}}
\toprule
\textbf{File} & \textbf{Role and important details} \\
\midrule
\endfirsthead
\toprule
\textbf{File} & \textbf{Role and important details} \\
\midrule
\endhead
\texttt{tests/closed\_loop.py} & Shared closed-loop simulation harness used by controller benchmark tests. \\
\texttt{tests/test\_hardware.py} & Unit tests for hardware abstractions (PCA9685 math, VESC protocol, backend logic). \\
\texttt{tests/test\_lookahead.py} & Tests for curvature-aware lookahead scheduling correctness and compatibility. \\
\texttt{tests/test\_map\_controller.py} & Tests MAP controller LUT interpolation and steering behavior. \\
\texttt{tests/test\_mpc.py} & Closed-loop MPC behavior and solve-budget validation tests. \\
\texttt{tests/test\_mpcc.py} & MPCC closed-loop corridor and tracking tests. \\
\texttt{tests/test\_race\_brain.py} & Tests opponent detection/fusion/strategy logic and race-mode transitions. \\
\texttt{tests/test\_racing\_tech.py} & Consolidated tests for MAP and velocity-profile race-tech modules. \\
\texttt{tests/test\_refiner.py} & Tests raceline-refinement optimizer logic. \\
\texttt{tests/test\_spliner.py} & Tests Frenet spliner geometry and overtake-path generation. \\
\texttt{tests/test\_state\_estimation.py} & Tests velocity EKF and lidar deskew behavior. \\
\texttt{test/test\_flake8.py} & Lint style gate for Flake8 (ament tooling). \\
\texttt{test/test\_pep257.py} & Docstring style gate for PEP257 (ament tooling). \\
\texttt{test/test\_copyright.py} & License/copyright header gate (ament tooling). \\
\bottomrule
\end{longtable}

\subsection{UI and packaging/support code}
\begin{longtable}{p{0.30\linewidth}p{0.66\linewidth}}
\toprule
\textbf{File} & \textbf{Role and important details} \\
\midrule
\endfirsthead
\toprule
\textbf{File} & \textbf{Role and important details} \\
\midrule
\endhead
\texttt{ui/server.py} & Lightweight dashboard backend; exposes REST endpoints, launches/stops race demos in the simulation container (Docker in this repository setup), and serves telemetry/raceline/image data. \\
\texttt{ui/index.html} & Browser dashboard frontend for tuning, raceline generation, and race monitoring actions. \\
\texttt{setup.py} & Python package setup and ROS 2 console entry points for executable nodes. \\
\texttt{setup.cfg} & Packaging/lint configuration metadata for Python project tooling. \\
\texttt{package.xml} & ROS package metadata and dependency declarations for build/runtime integration. \\
\texttt{tools/finish\_mapping.sh} & Shell workflow helper to finalize a mapping session and persist map/raceline artifacts. \\
\texttt{launch/ego\_racecar.xacro} & Robot description (URDF/Xacro) for ego vehicle visualization/transforms in sim. \\
\texttt{launch/opp\_racecar.xacro} & Robot description (URDF/Xacro) for opponent vehicle visualization/transforms in sim. \\
\bottomrule
\end{longtable}

\section{How major subsystems fit together}
\subsection{Mapping and raceline generation phase}
\begin{enumerate}
  \item Run \texttt{slam\_mapping\_launch.py} to start SLAM.
  \item Drive manually or with \texttt{mapping\_driver.py} to map the track.
  \item \texttt{track\_learner.py} monitors map growth and repeatedly calls \texttt{raceline\_optimizer.py}.
  \item Best feasible line is copied to \texttt{racelines/best\_raceline.csv}.
\end{enumerate}

\subsection{Race-time control phase}
\begin{enumerate}
  \item Load raceline CSV (x, y, heading, curvature, speed).
  \item Optionally refine curvature and recompute speed profile.
  \item Track line using MPC each tick.
  \item If solver unavailable/fails, execute MAP fallback command.
  \item Apply traction and emergency-braking safety overrides.
  \item Send actuation through \texttt{drive\_node} backend.
\end{enumerate}

\subsection{Perception and opponents}
\begin{itemize}
  \item Lidar-only opponent extraction and tracking: \texttt{race\_brain.py}.
  \item Camera-based YOLO opponent detections: \texttt{camera\_perception.py} (and optional OAK-D on-device detections in \texttt{oakd\_camera.py}).
  \item Fusion and strategic behavior (attack/defend/evade): \texttt{race\_agent.py} + \texttt{race\_brain.py}.
\end{itemize}

\section{What is actually deployed on the vehicle for a robust race setup}
For the cleanest competition stack in this repository, the deployed nodes are typically:
\begin{itemize}
  \item \texttt{rplidar\_node}
  \item \texttt{oakd\_camera}
  \item \texttt{velocity\_ekf}
  \item localization node (external, typically particle filter)
  \item \texttt{raceline\_mpc}
  \item \texttt{drive\_node}
\end{itemize}

Among these, \texttt{raceline\_mpc.py} is the primary race logic loop and \texttt{drive\_node.py} is the hardware command execution/safety gate.

\section{Practical reading order for newcomers}
\begin{enumerate}
  \item \texttt{README.md} for architecture and startup commands.
  \item \texttt{launch/car\_bringup\_launch.py} and \texttt{config/hardware.yaml} to see real-car wiring.
  \item \texttt{f1tenth\_gym\_ros/raceline\_mpc.py} and \texttt{f1tenth\_gym\_ros/mpc\_controller.py} for main control logic.
  \item \texttt{f1tenth\_gym\_ros/drive\_node.py} for actuation/safety.
  \item \texttt{f1tenth\_gym\_ros/raceline\_optimizer.py} for map-to-raceline generation.
  \item \texttt{tests/} for behavior guarantees and expected invariants.
\end{enumerate}

\section{Final summary}
AtlasAutoware is a complete autonomous racing stack where:
\begin{itemize}
  \item map creation and raceline optimization happen first,
  \item race-time control is centered on \texttt{raceline\_mpc} at 50 Hz,
  \item \texttt{drive\_node} safely turns commands into motor/steering actions,
  \item and optional perception/strategy modules support multi-car racing.
\end{itemize}

This document is intended to give a full beginner-to-codebase bridge: terminology, architecture, runtime behavior, and per-file responsibilities.

\end{document}
